\chapter{The Thorin measure at level one and the backward Cauchy flow}\label{chap:lo}\chaptermark{Level one and the backward flow}

Chapter~\ref{chap:lv} shows that for $\theta\ge1$ the function $\xi(\theta)/\xi(\theta+\sqrt s)$ is, unconditionally, the Laplace transform of a generalised gamma convolution, that its Thorin measure is the Cauchy smoothing of the zero measure at width $\theta-\Re\rho$, and that the family $\theta\mapsto U_\theta$ is a Cauchy semigroup above the zero free abscissa (Theorems~\ref{lv:thm:thorin} and~\ref{lv:thm:semigroup}). It asks whether the semigroup can be run backward from a level one measure computed without the zeros (Problem~\ref{lv:prob:down}). This chapter does three things with that question.

First, it computes the level one measure in closed form. Its density is $\pi^{-1}\Re(\xi'/\xi)(1+iv)$, a smooth gamma part minus the prime sum $\pi^{-1}\sum_{n\ge2}\Lambda(n)n^{-1}\cos(v\log n)$, read as a compensated limit because the series itself diverges. The Fourier transform of the measure is, apart from the transform of the smooth part and a constant $1$ from the pole of $\zeta$, a sum of point masses $-\Lambda(n)/n$ at $\omega=\pm\log n$: the level one measure carries the primes in its spectrum and the zeros as the centres of its Cauchy components, and the identity between the two descriptions is a Weil type explicit formula for one analytic test function evaluated at the zeros.

Second, it studies the backward flow. On the Fourier side the semigroup multiplies by $e^{-(\theta-\theta')|\omega|}$, so going from level $1$ to level $\theta'$ multiplies by $e^{(1-\theta')|\omega|}$. The band limited backward flow converges to the level $\theta'$ measure if and only if $\zeta$ has no zeros to the right of $\theta'$; a zero to the right makes the limit fail to be a measure $\mu$ with $\int|\mu|(dv)/(1+v^2)<\infty$. Numerically, from samples of the level one density on the line $\Re s=1$ only, the flow places a peak within $0.035$ of each of the $109$ zeros tested, in $100<v<300$.

Third, it treats the backward flow as a positive linear inverse problem of Vardi and Lee type with the Cauchy kernel as the blur. The EM algorithm of Shepp and Vardi, an information geometric analogue of von Neumann's alternating projection (Csisz\'ar and Tusn\'ady), gives a positivity preserving backward flow with a diagnostic, the total Kullback Leibler residual; its data augmentation Gibbs sampler is its Monte Carlo version; and generative Bayesian computation amortises a posterior under a GUE prior, trained on simulations. The EM inversion locates the $109$ zeros without ringing, and a redesigned injection test separates an injected off line quartet from a control; the GBC implementation, rerun here with a corrected forward model, counts the zeros in unit bins almost exactly.

The chapter ends by asking whether any of this is a path to the Riemann hypothesis. It is not. It is a dictionary in which RH reads as the statement that the backward flow runs all the way to $\theta=\tfrac12$, and in which the exponential amplification $e^{(1-\theta')|\omega|}$ is the classical error term $\psi(x)-x=O(x^{\theta'+\varepsilon})$ in Fourier clothing.

The chapter develops a source chapter by Nick Polson \cite{P26lo}, together with the computations done for this book (the scripts \texttt{code/lo\_common.py}, \texttt{code/lo\_flow.py}, \texttt{code/lo\_em.py} and \texttt{code/lo\_gbc.py}, their outputs \texttt{data/lo\_*.json} and logs, listed in \texttt{code/README.md}; the scripts of the source were not used, and the sampler of $u_1$ and the EM iteration are those of Chapter~\ref{chap:cx}, \texttt{code/cx\_em.py}). Every number printed below is one of these recomputations unless it is marked as the author's or the review's, and where the recomputation differed from the source the recomputed value is printed. Claims of the source that could not be reproduced are dropped or restated as such in the text.

Notation. The book's $\Re$, $\R$ and the cut plane $\cut$ are used; $\xi(s)=\tfrac12s(s-1)\pi^{-s/2}\Gamma(s/2)\zeta(s)$, a zero is $\rho=\sigma_\rho+i\gamma_\rho$, and $P_c(v)=(c/\pi)/(c^2+v^2)$ is the Cauchy kernel of width $c>0$, with $P_0=\delta_0$ (the Poisson kernel of Section~\ref{ch:clock}, with the same index). The Thorin measure of this chapter is written $\mathfrak U_\theta$ (Section~\ref{lo:sec:level}). It carries no factor $\tfrac12$ and is therefore twice the symmetric measure $\widetilde U_\theta$ of Theorem~\ref{lv:thm:semigroup}; its density is the density $\rho_{\theta-1/2}$ of Chapter~\ref{chap:tf} (Theorem~\ref{tf:thm:deconv}). The digamma function is written $\Gamma'/\Gamma$, because $\psi(x)=\sum_{n\le x}\Lambda(n)$ is Chebyshev's function here; the grid spacing of the numerics is $h$ and the test function of the explicit formula is $\varphi$, the Fourier transform being $\widehat f(\omega)=\int e^{-i\omega v}f(v)\,dv$.

Status words follow the rest of the book. A statement is proved when a proof or a pointer is given, computed when a script of the book produced it, conditional, heuristic or open as stated; nothing in this chapter is certified in ball arithmetic except the sampler of $u_1$, which is evaluated with python-flint balls at $128$ bits and checked at $192$. The results are these.
\begin{enumerate}
\item The level one Thorin measure of $\xi$ has the explicit real analytic density \eqref{lo:eq:u1}, for every real $v$ including $v=0$, where the prime sum is a compensated limit with the value $u_1(0)=(2+\gamma_E-\log4\pi)/(2\pi)$ (Theorem~\ref{lo:thm:density}, proved; for $v\ne0$ Propositions~\ref{bp:prop:thorin} and~\ref{bp:prop:sin2}). Its positivity is automatic from the Euler product; the nonvanishing of $\zeta$ on $\Re s=1$ is equivalent to the absence of atoms (Section~\ref{lo:sec:path}, proved).
\item The Fourier transform of the measure is the sum of prime atoms, the transform of the smooth part and the constant $1$, and the identity of the two descriptions is \eqref{lo:eq:weil}, an explicit formula for one analytic test function (Proposition~\ref{lo:prop:spectrum}, proved). The spectral peaks lie at $\log n$ to $10^{-4}$ (Table~\ref{lo:tab:spectrum}, computed).
\item The band limited backward flow converges at level $\theta'$ if and only if $\zeta$ has no zeros in $\Re s>\theta'$, and RH is the convergence at every $\theta'>\tfrac12$ (Theorem~\ref{lo:thm:dichotomy}, proved). The flow recovers $109$ zeros to $0.035$ at $\theta'=\tfrac34$ and to $0.027$ at $\theta'=\tfrac12$; the sharp cutoff rings and the Fej\'er window does not (Section~\ref{lo:sec:numflow}, computed).
\item The existence of a positive solution of the deconvolution $P_{1-\theta'}*x=u_1$ is the zero free half plane $\Re s>\theta'$, and the total Kullback Leibler residual is bounded away from zero as the grid is refined exactly when a zero lies to the right of $\theta'$ (Proposition~\ref{lo:prop:emres}, proved, with a counterexample to the residual per unit length). The EM runs and a quartet injection test with controls are computed (Section~\ref{lo:sec:emnum}).
\item The data augmentation sampler needs an exposure and pseudo counts; it is a method and carries no result. Generative Bayesian computation under a GUE prior counts the zeros in unit bins with $98$ per cent exact bins on the $200$ test bins, for a corrected forward model (Section~\ref{lo:sec:gbc}, computed; the formulation is heuristic).
\item None of this is a path to RH (Section~\ref{lo:sec:path}); four open problems close the chapter (Section~\ref{lo:sec:open}).
\end{enumerate}
